\section{The Loewner family and the flipped-return loop}
\label{sec:loewner}

\subsection{Scaling collapse}

\begin{proposition}[conformal collapse]
\label{prop:collapse}
For the linear driver $u(t)=at$ the trace satisfies $q_a(Tr)=\sqrt T\,q_{a\sqrt T}(r)$. In \Nfour{} SYM,
dilatation invariance of renormalized loops, up to cusp anomalous dimensions, gives
\begin{equation}
W(a,t)=F\big(a\sqrt t,\ \delta/\sqrt t\big)\,(\mu\sqrt t)^{-\Gamma_{\rm tot}(a\sqrt t)},
\end{equation}
where $\delta$ is a rounding scale and $\Gamma_{\rm tot}$ the sum of the cusp anomalous dimensions.
\end{proposition}

The Loewner time is therefore not an independent clock in the conformal theory: renormalized loops depend
on $t$ only through the slope $s=a\sqrt t$, and the short-time limit is the small-driver limit. Pure
Yang--Mills at fixed flow resolution differs, because there $t/\tau$ is meaningful.

\subsection{Choice of completion}

The trace must be closed to form a loop. Three completions were examined in units where $t=1$, where the
trace and a straight return chord bound a sliver of width $w(y)=(s/6)y(2-y)$, $y\in[0,2]$, with cusp openings
$s/3$ at base and tip.

\emph{Constant scalar on both legs.} The opposite legs are nearly antiparallel with the same scalar charge, so
gauge and scalar exchange add. At one loop $\log W\approx(3\lambda(1-1/N^2)/2\pi s)\log(1/\hat\varepsilon)$,
which diverges as $s\to0$; to all orders the near-antiparallel cusp gives $\log W\approx(6c/s)\log(1/\hat\varepsilon)$
with $c$ the antiparallel-lines coefficient \cite{DrukkerForini}. The $a=0$ member is not a normalizable
control. At fixed flow resolution the loop measures perimeter growth, $W_\tau=1+8t\vev{N^{-1}\tr\Phi_{n_0}^2}_\tau+\dots$.

\emph{Tangent-coupled scalar.} With $n=MT$ the loop is a Zarembo loop and $\vev W=1$ for every curve
\cite{Zarembo,GuralnikKulik,DGGM}; the corners are BPS cusps. This is a Ward-identity control, reproduced by
Proposition~\ref{prop:quadratic} at $O(\epsilon^2)$ exactly because $C_D=6C_\Phi$.

\emph{Flipped return.} Keep $n_0$ on the trace and $-n_0$ on the chord. The chord traversed from tip to base
with connection $\ii A\cdot\dd x-|\dd x|\Phi$ is the inverse of the base-to-tip transport $R_t$ with connection
$\ii A\cdot\dd x+|\dd x|\Phi$, so
\begin{equation}
Q_t=R_t(1)^{-1}U_t,\qquad W(t)=\vev{N^{-1}\tr Q_t},
\end{equation}
and at $a=0$ the transport is exactly the identity. The opposite legs are nearly antiparallel with opposite
scalar charge (the no-force configuration), and both cusps have $(\varphi,\theta)\approx(\pi-s/3,\pi)$, i.e.\ they
are nearly BPS. $W$ is even in $a$. This is the recommended family.

\subsection{Exact evolution equations}

Write $k=q_1\dot q_2-q_2\dot q_1$, $\hat n=(-q_2,q_1)/|q|$, $\rho=|q|$, $\gamma=|\dot q|-\dot\rho\ge0$; hats denote
conjugation by $R_t(r)$; $J(r)=\int_0^rr'\hat F_{12}\dd r'$, $\mathcal G(r)=\int_0^rr'\widehat{(\hat n\cdot D\Phi)}\dd r'$,
$J=J(1)$, $\mathcal G=\mathcal G(1)$.

\begin{proposition}[first equation]
\label{prop:first}
For smooth fields and every configuration,
$\dot Q_t=\mathcal X_tQ_t$ with $\mathcal X_t=\ii k(J+\ii\mathcal G)+\gamma\hat\Phi(1)$. Since
$\int_0^t\gamma=L_{\rm trace}(t)-|q(t)|$, the scalar tip coefficient is exactly the growth rate of the perimeter
mismatch between the two legs.
\end{proposition}

\begin{proposition}[transported derivative and second equation]
\label{prop:second}
For a local adjoint field $X$ on the chord,
$\partial_t\hat X(r)=r\widehat{(\dot q\cdot DX)}(r)+[\hat X(r),\mathcal K(r)]+\widehat{(\partial_tX)}_{\rm explicit}$
with $\mathcal K(r)=-\ii k(J(r)+\ii\mathcal G(r))+\dot\rho\,r\hat\Phi(r)$. Hence
$\ddot W=\vev{N^{-1}\tr((\dot{\mathcal X}+\mathcal X^2)Q)}$ with
\begin{align}
\dot{\mathcal X}&=\ii\dot k(J+\ii\mathcal G)+\dot\gamma\hat\Phi(1)+\ii k(\dot J+\ii\dot{\mathcal G})
+\gamma\big[\widehat{(\dot q\cdot D\Phi)}(1)+[\hat\Phi(1),\mathcal K(1)]\big],\\
\dot J&=\int_0^1r\big(r\widehat{(\dot q\cdot DF_{12})}+[\hat F_{12},\mathcal K(r)]\big)\dd r,\\
\dot{\mathcal G}&=\int_0^1r\Big(r\widehat{(\dot q^\mu\hat n^\nu D_\mu D_\nu\Phi)}-\frac{k}{|q|^2}\widehat{(\hat q\cdot D\Phi)}
+[\widehat{(\hat n\cdot D\Phi)},\mathcal K(r)]\Big)\dd r .
\end{align}
Both equations hold for any driver with a $C^2$ trace.
\end{proposition}

The proofs use the ordered-transport variation on the chord and one radial integration by parts, which is
possible because the chord connection commutes with $\Phi$ along the chord; the commutator terms arise
because it does not commute with $F$ or $D\Phi$. Both equations were verified by finite differences on a
noncommuting $SU(2)$ background (relative errors $10^{-9}$ and $8\times10^{-7}$). The terms
$\ii k\,\ii k[\hat F,J(r)]$ and $(\ii kJ)^2$ combine into $-2k^2\mathcal P$ with $\mathcal P$ the ordered
$\hat F\hat F$ moment, exactly as in pure Yang--Mills; the new terms are ordered pairs mixing $\hat F$,
$\widehat{(\hat n\cdot D\Phi)}$ and $\hat\Phi$ and chord--tip pairs, which do not combine into a single moment.

\subsection{Short time at fixed resolution}

Configuration by configuration, with $\mathcal A=-(2a/9)t^{3/2}$ the oriented area and
$L_{\rm trace}-L_{\rm chord}=a^2t^{3/2}/27$,
\begin{equation}
N^{-1}\tr Q_t-1=\frac1{2N}\tr\Big(\mathcal A(\ii F_{12}+D_1\Phi)(0)+\frac{a^2t^{3/2}}{27}\Phi(0)\Big)^2+O(t^{7/2}),
\end{equation}
so that in an $SO(6)$- and parity-invariant flowed state
\begin{equation}
W_\tau=1+\frac{2a^2t^3}{81}\big(C^\Phi_\tau-C_\tau\big)+\frac{a^4t^3}{1458}S_\tau+O(t^{7/2}),\qquad
\log W^{(1)}_\tau=B_1\Big[-\frac{a^2t^3}{1296\tau^2}+\frac{a^4t^3}{5832\tau}\Big]
\end{equation}
with the free flowed propagator ($C_\tau=\vev{N^{-1}\tr F_{12}^2}_\tau$, $C^\Phi_\tau=\vev{N^{-1}\tr(e\cdot D\Phi_{n_0})^2}_\tau$,
$S_\tau=\vev{N^{-1}\tr\Phi_{n_0}^2}_\tau$), and an independent
expansion of the flowed one-loop Maldacena--Wilson integral for $\tau\gg t$ reproduces it term by term through the
exact identities $M_0=(L_{\rm trace}-L_{\rm chord})^2$ and $M_2=2\mathcal A^2$. The agreement is not circular:
one side is the exact first equation, the other a Feynman integral.

\subsection{Continuum one loop and the flow regimes}

With cusp neighbourhoods of size $\varepsilon$ excised,
\begin{equation}
\begin{gathered}
\log W=B_1\Big[-A(s)\log\frac\ell\varepsilon+F(s)\Big]+O(\varepsilon),\\
A(s)=\sum_{\rm cusps}2(\pi-\alpha_c)\tan\frac{\alpha_c}2=\frac{2\pi}3s+A_2s^2+\dots,\qquad F(s)=\frac{4\pi}3s+O(s^2\log\tfrac1s),
\end{gathered}
\label{eq:oneloop}
\end{equation}
$A_2\approx-0.224$. The exactness of $A$ is tested without fitting: after subtracting it, the successive
$\varepsilon$-differences halve. So $W\to1$ like $1-(2\pi B_1/3)|s|(\log(\ell/\varepsilon)-2)$, non-analytic in $a$
because $a\to-a$ is a symmetry. With a gradient-flow regulator of resolution $\tau$ there are three regimes:
regime I ($\tau\gg t$) is the $a^2t^3$ law above; regime II ($s^2t\ll\tau\ll t$) gives
$\log W\approx-B_1(\sqrt{2\pi}/27)s^2\sqrt{t/\tau}$; regime III ($\tau\ll s^2t$) gives
$\log W\approx-B_1(2\pi s/3)[\log(8\kappa)+\gamma_E/2-2]$ with $\kappa=s\sqrt t/(6\sqrt{8\tau})$. Separation-based
regulators cut the near-BPS cusp interaction where the sliver width reaches $\sqrt\tau$, at distance
$\approx\sqrt\tau/\alpha$ from the vertex, which produces an $|s|\log|s|$ term absent from vertex excision; the
scheme-independent content at $O(s)$ is $A(s)$. Every regime is regular, and the order of the limits $t\to0$ and
$\tau\to0$ changes the exponent but never produces a divergence.

\subsection{Rounded family and closure}

Replacing each cusp by a circular fillet of radius $\delta$, with the internal vector turning in step with the
tangent, gives a $C^1$ loop with a finite one-loop value
\begin{equation}
\log W^{(1)}=-B_1\Big[A(s)\log\frac1{\hat\delta}+\frac{2\pi|s|}{3}\Big(\log\frac{|s|}3-2\Big)-h|s|\Big]+O(s^2),\qquad
\hat\delta=\delta/\sqrt t,\quad h=-4.0202,
\end{equation}
valid for $\hat\delta\ll|s|$. The coefficient of $\log(1/\hat\delta)$ is the exact cusp sum and the $|s|\log|s|$ term
has the same origin as in regime III; the constant $h$ is convention dependent (an independent implementation gave
$-4.02025$).

The second equation is not closed. At a finite regulator the Schwinger--Dyson identities reduce only the
curvature-current part $(\dot q^1E_2-\dot q^2E_1)$, the scalar Laplacian and the radial derivatives, leaving
the transverse gauge currents, the transverse scalar Laplacian $(D_3^2+D_4^2)\Phi$, scalar and fermion currents,
scalar-potential and Yukawa insertions, scalar commutators from transport along the chord, the tip gradient and
the ordered pairs. At one loop the field-equation part of a derivative moment is a pure regulator contact term
while the transverse part has a long-range $4/r^4$ kernel, so the residual cannot be dropped even
perturbatively. The flipped hierarchy is an honest infinite hierarchy with a specific residual inventory.

\subsection*{Ledger}
\begin{center}\small
\begin{tabular}{@{}p{0.55\textwidth}p{0.40\textwidth}@{}}
\toprule
Statement & Status\\
\midrule
Conformal collapse, Proposition~\ref{prop:collapse} & Proof given conformal invariance\\
Constant-scalar completion diverges as $1/s$ & One loop: labelled computation; all orders: proof sketch\\
Zarembo family trivial & Literature\\
Identity at $a=0$; evenness in $a$ & Proof\\
Propositions~\ref{prop:first}, \ref{prop:second} & Proof (smooth fields); numerically verified\\
Short-time formula and one-loop confirmation & Proof sketch (remainder not bounded); independent one-loop check\\
Continuum $A(s)$; $F(s)$ leading term & Proof (cusp formula); exactness tested; $s^2\log$ term an observation\\
Regimes I--III & Analytic (thin sliver, moments) plus numerics\\
Rounded family; $h$ & Numerical verification; $h$ a labelled constant\\
Closure inventory & Structural reading with one-loop kernel controls\\
\bottomrule
\end{tabular}
\end{center}
